% Section 2 (2026-10-27). Every factual sentence ends with a comment naming its ledger rows (ledger-v11.tsv).
\section{Background}\label{sec:background}

\subsection{Comparator and nanoda}

Comparator's README at commit d03acab \cite{comparator} states three guarantees for each listed theorem of a solution: that it proves the same statement as provided in the challenge, that it uses no more axioms than listed in permitted\_axioms, and that it is accepted by the Lean kernel. % v11:L102 v11:L103 v11:L104
The same README states that external kernels registered in external\_kernels are given a file containing the solution export. % v11:L104
According to Comparator's and lean4export's READMEs at the pinned commits, both kernel checks depend on the Lake build of the solution in a landrun sandbox, on the lean4export export of the solution, and on Comparator's own statement and axiom checks. % v11:L142
By the same READMEs, only the Lean kernel's check replays the solution environment inside the Comparator process, and only nanoda's check runs as a separate process given a file containing the solution export and a nanoda-style config.json. % v11:L142
Comparator's README states the guarantees under six assumptions, among them that ``The Lean kernel is correct (with `external\_kernels` this can be reduced to "At least one of the Lean kernel or the `external\_kernels` is correct")''. % v11:L137
nanoda's README at commit 3a24072 \cite{nanoda-lib} states that nanoda\_bin takes one argument, a JSON configuration file, and reads an export file created by lean4export. % v11:L105

All runs used Lean 4.34.1 (5045d005), Mathlib d13f23b, Comparator v4.34.0 d03acab and lean4export v4.34.0 076e8e5 \cite{lean4,mathlib,comparator,lean4export}. % v11:L087
The full commits are d03acab154d269c06e60e4de7e4cc85deebff94b for Comparator and 076e8e57707e813375e8f9da8bf989799ace9680 for lean4export. % v11:L092
The official nanoda is 3a24072; f67473b, 37cb87b and 37cb87b with nealmcb's patch at 64 MiB are unofficial builds, reported separately. % v11:L088

\subsection{The challenge}

The SiegelZeros challenge file has sha256 bedaac0248a641910f79a380d9702d82c6bfbbfd3760a69d27c72c5d8fb136ce, and it is unchanged between adc7f12 and fd4aeeb. % v11:L007
Its config has Comparator check two theorems, OAI.SiegelZeros.WeightedTorusJets.exists\_absolute\_real\_zero\_gap and OAI.SiegelZeros.WeightedTorusJets.dirichletRealZeroBound\_proof. % v11:L147
Section~\ref{sec:faithfulness} quotes both statements, and Section~\ref{sec:gap} gives the Lean proof that they are equivalent. % v11:L147 v11:L251

\subsection{The paper's lemmas}

The paper \cite{openai-siegel-2026} states its Theorem 1 at refs/siegel-paper.tex:66-73, Lemma 2 at 181-194, Lemma 3 at 252-267, Corollary 4 at 410-420, Lemma 6 at 517-524 and Lemma 7 at 598-604. % v11:L106 v11:L107 v11:L108 v11:L109 v11:L110 v11:L111
The paper numbers all of these with one counter, so the result between Corollary 4 and Lemma 6 is Remark 5. % v11:L106 v11:L110
Lemma 2 gives a prime-sum bound for X >= 3 and a lower bound for the mass of primes with chi(p) = -1. % v11:L107
Lemma 3 states that, for a surjective linear map A from C\textasciicircum4 to C\textasciicircum3 whose kernel is spanned by a vector with Q-linearly independent coordinates, evaluation at A n for n in E\_N is surjective from P\_t under the budget (3.1). % v11:L108
The paper cites Fischler \cite{fischler2005} as [Fischler, Theorem 1.1 and Lemma 1.3] at refs/siegel-paper.tex:155. % v14:L273
Corollary 4 states that, under the hypotheses of Lemma 3, the monomial rows with alpha\_j <= T\_j span C\textasciicircum(E\_N). % v11:L109
Lemma 6 bounds the exponent sums S\_1 and S\_2, and Lemma 7 states the divisibility of Delta by p\textasciicircum(E\_p) for every admissible prime p. % v11:L110 v11:L111
